\documentclass[journal]{IEEEtran}

\usepackage{amsmath,amssymb}
\usepackage{booktabs}
\usepackage{graphicx}
\usepackage{hyperref}
\usepackage{xcolor}
\usepackage{placeins}
\usepackage{balance}
\hypersetup{hidelinks}

\definecolor{draftred}{RGB}{150,20,20}
\newcommand{\draftwarning}{\textcolor{draftred}{DRAFT---NOT FOR SUBMISSION}}

\title{Explicit Process Relations Prevent Semantic Misrouting in Long-Horizon
Robotic Wire-Harness Assembly}
\author{Anonymous Authors}
\date{}

\begin{document}
\maketitle

\begin{center}
\draftwarning
\end{center}

\begin{abstract}
Automotive wire routing requires geometric control and process semantics:
cable segments must enter designated fixtures in order. We isolate semantic
misrouting through paired interventions that preserve geometry, layout, physics,
and seed while changing
only the feasible segment--fixture assignment. Our controller represents cable
connectivity and process correspondence as typed graph relations and grounds
actions with graph pointers. Across 300 paired PBD episodes, the relational
policy succeeds in 251 cases (83.7\%) versus 10 (3.3\%) for an equal-parameter
geometry model; paired advantages are 90, 83, and 68 points across three
difficulties. All relation-aware models share the same 251/300 outcome,
supporting relation availability rather than one graph architecture. In
nine-assignment counterfactuals, relation-equipped policies
succeed in 2,340/2,700 cells versus 57/2,700 for Geometry. A direct MuJoCo audit
gives 180/180 versus 0/180. An offline audit on 1,647 public real-robot
cable-routing trajectories confirms a temporal action signal but not semantic
success. DLO-Lab distinguishes ordered from wrong-order paths, yet a direct
endpoint path scores similarly to the ordered path, exposing limited semantic
discrimination. The evidence isolates the need for observable process
relations; trained external semantic replication and force-controlled hardware
remain open.
\end{abstract}

\begin{IEEEkeywords}
Deformable object manipulation, graph representations, industrial robotics,
semantic task planning, wire-harness assembly.
\end{IEEEkeywords}

\section{Introduction}

Deformable linear objects (DLOs) combine high-dimensional geometry with
contact-rich dynamics.  Industrial wire harnesses add a second difficulty:
the identities of endpoints, branches, terminals, and fixture-bound segments
matter to the manufacturing process.  Recent systems have demonstrated
contact-aware cable shaping~\cite{chen2023contact}, twisting-based harnessing
with force control~\cite{zhang2024twisting}, and integrated multi-robot DLO
assembly~\cite{chen2025multirobot}.  Concurrent benchmarks now cover
industrial clip routing and connector insertion~\cite{zhu2026wirecraft}, as
well as broad DLO manipulation with differentiable physics
~\cite{cao2026dlolab}.  Consequently, our contribution is not another claim of
being the first industrial DLO simulator.

We instead isolate a semantic misrouting problem.  Three cable particles may
be physically similar and spatially adjacent, but each is bound to a different
clip.  Once the wrong segment is latched, low geometric error is not enough to
recover the process.  This motivates an intervention in which only the
segment--clip relation changes.  Geometry-only shortcuts should be invariant to
that change, whereas a relational controller should change its decision while
preserving performance.

The intended contributions are:

\begin{enumerate}
  \item a paired nine-assignment intervention, frozen by an independent
  teacher-feasibility screen, that holds the physical RouteBot episode fixed
  and changes only the segment--clip process relation;
  \item a typed relational state and graph-pointer policy for long-horizon
  grasp, transport, latch, and endpoint-placement decisions;
  \item a controlled comparison with equal-parameter geometry ablation and an
independently implemented ACT-style state baseline with direct, typed-pointer,
and fixed relational-pooling variants;
  \item failure-aware statistics, an attributed public real-data audit, and
  cross-physics replication that distinguish semantic correctness from
  geometric completion.
\end{enumerate}

\section{Related Work}

\subsection{Industrial DLO assembly}

Contact-aware planning and object-centered skills can maintain DLO shape using
fixtures and force feedback~\cite{chen2023contact}.  Twisting-induced tension
enables a single arm to insert wires into clamps in narrow workspaces
~\cite{zhang2024twisting}.  Multi-robot pipelines combine perception, handover,
and placement across the assembly workflow~\cite{chen2025multirobot}.  These
systems provide stronger real-physics evidence than our present simulation and
bound our claims to semantic decision making.

WireCraft spans connector insertion, clip routing, and channel seating with
articulated and deformable physics~\cite{zhu2026wirecraft}.  DLO-Lab emphasizes
material diversity, differentiable physics, grasp sensitivity, and long-horizon
decomposition~\cite{cao2026dlolab}.  Our experiments target a narrower but
orthogonal question: whether a policy respects process relations under exact
semantic counterfactuals.  External replication on one of these benchmarks is
a required final experiment.

Luo et al. collected real Franka Panda demonstrations for single- and
multi-clip cable routing and showed that learned recovery and primitive
selection matter over long horizons~\cite{luo2024multistage}.  We use its public
LeRobot conversion only as an episode-disjoint offline action-interface audit;
the release does not expose our segment--clip relation intervention.

\subsection{Learning DLO dynamics and action sequences}

Graph networks have modeled cable deformation for offline--online learning and
MPC~\cite{wang2022offlineonline}; graph structure alone is therefore not a new
contribution.  Our graph edges instead combine physical connectivity with a
manufacturing correspondence relation.  ACT predicts action chunks with a
conditional variational Transformer~\cite{zhao2023act}, while Diffusion Policy
models multimodal action sequences through conditional denoising
~\cite{chi2023diffusion}.  We use a state-based ACT adaptation as the matched
temporal baseline and do not claim a comprehensive comparison of every
visuomotor sequence model.

\section{Problem Formulation}

Let $G_t=(V_t,E_t^{p},E_t^{s})$ denote the observable state at time $t$.
$E_t^{p}$ contains cable connectivity; $E_t^{s}$ contains process relations
between cable segments, clips, the finish fixture, and the end effector.  A
route is successful only if every clip $c_i$ is bound to its assigned particle
$\pi(i)$ and the endpoint reaches its fixture.  Because a crossing in a 2-D
projection need not be a topological crossing of the real 3-D harness, we
report crossing-free completion as a stricter secondary safety metric rather
than mixing it into semantic success. We report geometric coverage separately because
$\lvert\{c_i\ \text{latched}\}\rvert=3$ does not imply semantic correctness.

For a physical episode seed $s$, the counterfactual family selects one cable
particle from each of three non-overlapping local arc-length windows.  Before
the final learned-policy evaluation, we screened nine order-preserving candidates
on ten held-out seeds per difficulty. All 270 semantic-completion rollouts
succeeded, so we froze all nine assignments. Cable coordinates, clip coordinates, obstacle geometry,
material parameters, integration noise, and initial arm state are identical.
Only $E_t^{s}$ and the corresponding success predicate change.  This makes each
policy comparison paired at the physical-seed level.

\section{Method}

\subsection{Typed relational observation}

Each node contains normalized position, generic role, endpoint/junction role,
process order, activation state, confidence, progress, and difficulty.  The
policy adds self-loops and performs message passing over
$\mathrm{clip}(E^p+E^s,0,1)$.  Attention pooling is concatenated with global
state and a task embedding.

\subsection{Graph-pointer action grounding}

Two outputs are learned jointly: continuous action logits and pointer logits
over nodes.  RouteBot uses one pointer.  When no segment is held, the decoder
restricts candidates to unbound cable particles and closes the gripper only
inside a fixed observable distance.  While carrying a segment, candidates are
restricted to inactive fixture/finish nodes.  The policy never reads hidden
success labels.

\subsection{Baselines}

The no-topology ablation keeps model parameter count and training schedule
fixed while removing adjacency, endpoint/junction identity, process order, and
semantic identity.  The ACT-style baseline uses a CVAE Transformer, eight-step
action chunks, temporal ensembling, and the same typed node observations but no
graph edges.  We evaluate direct continuous ACT and ACT-Pointer, which shares
the candidate filter and learns a node-classification head.  The latter
controls for the discrete action interface. ACT-RelPool additionally receives
the same semantic adjacency but replaces learned graph propagation with fixed
one-hop neighbor averaging before its Transformer, controlling for access to
the relation itself.

\section{Experimental Protocol}

The fast backend uses position-based dynamics~\cite{muller2007pbd} for paired
large-sample experiments.  Each main-table cell contains at least 100 unseen
seeds at difficulties $0.2$, $0.5$, and $0.8$.  All policies receive identical
seeds.  Binary success uses Wilson intervals; paired differences use McNemar's
exact test and a matched odds ratio.  We additionally report clip coverage,
topology accuracy, wrong latches, endpoint error, crossings, stretch error,
steps, and a mutually exclusive failure taxonomy.

The frozen imitation corpus contains 432 complete teacher trajectories and
19,190 state--action samples (SHA-256 prefix \texttt{dac7dd2b8ac}).  Thirty
training, nine validation, and nine test physical scenes are mutually disjoint.
Every physical scene is replayed under all nine assignments, yielding 270/81/81
training/validation/test trajectories; thus assignment identity cannot be
confounded with geometry within a scene.  The target-particle identity is not a
unary feature: changing an assignment changes only semantic adjacency edges and
the success predicate.

TopoHarness and its Geometry ablation use identical 64-dimensional, three-layer
networks and optimization schedules.  Both ACT variants use a 48-dimensional
one-layer CVAE Transformer, eight-action chunks, and a two-simulator-step control
stride matching demonstration sampling.  Model selection uses validation loss;
no frozen evaluation seed is used for architecture or assignment selection.

For external real-data grounding, we pin the public Berkeley Cable Routing
LeRobot conversion at commit
\texttt{20a7774a714a5daf56b220e01635425f2cb9745b}.  Its 42,328 frames form
1,647 complete real-robot teleoperation episodes with four camera views.  We
split whole episodes 70/15/15 into 1,153/247/247 train/validation/test episodes,
tune ridge regularization on validation episodes only, and report the four
nonconstant Cartesian action dimensions.  A state-only ridge, previous-action
control, and state-plus-previous-action ridge are evaluated on the same held-out
episodes.  The latter is an audit baseline, not our semantic-routing method.

\begin{figure}[t]
  \centering
  \includegraphics[width=\columnwidth]{%
    ../../artifacts/papers/routebot/public_real/berkeley_real_multiview.png}
  \caption{Aligned views from the pinned public real-robot cable-routing data
  (Berkeley Cable Routing, CC BY 4.0).  Images are evidence of the external data
  interface, not RouteBot rollouts.}
  \label{fig:berkeley-real}
\end{figure}

The three-dimensional replication uses MuJoCo~\cite{todorov2012mujoco} cable
elasticity.  Detailed UR10e and Robotiq meshes are kinematic visual twins with
collisions disabled; they must not be interpreted as torque-controlled robot
dynamics.  For the direct audit, each control step reconstructs the same
29-node typed graph from current MuJoCo cable, fixture, finish, and task-space
gripper poses.  The frozen network output is mapped back to a MuJoCo mocap
target; grasp, latch, and endpoint events are detected solely from native
MuJoCo body distances and constraints.  No PBD state or event is read.  Small
paired fixture perturbations define physical seeds, while the nine-assignment
intervention changes only semantic graph edges and the success predicate.

For an independently maintained simulator sanity check, we pin DLO-Lab at
commit \texttt{c5026a9} and its prescribed Mushroom-RL fork at
\texttt{ec33647}.  In the official \texttt{wiring\_post} environment, four
deterministic open-loop paths run in one matched four-environment batch: no
action, a straight endpoint move, a deliberately wrong fixture order, and the
target-consistent fixture order.  We report DLO-Lab's native geometric reward
without modification, plus ordered curve and fixture-relation diagnostics as
an explicitly labeled RouteBot extension.  This experiment tests executable
cross-simulator plumbing and metric sensitivity, not a trained-policy claim.

\section{Results}

\begin{table}[t]
  \centering
  \caption{Paired RouteBot results over all frozen difficulties. Semantic
  success is primary; crossing-free success is the strict 2-D shape metric.}
  \label{tab:main}
  \IfFileExists{../../artifacts/papers/routebot/final/main_table/route_main_table.tex}
    {{\small\setlength{\tabcolsep}{2.5pt}%
      \input{../../artifacts/papers/routebot/final/main_table/route_main_table.tex}}}
    {\fbox{\parbox{\dimexpr\columnwidth-2\fboxsep-2\fboxrule\relax}{%
      \centering Frozen table generated by
      \texttt{sim-eval-route-paper}}}}
\end{table}

\paragraph{Frozen paired PBD main table.}
TopoHarness completes 251/300 episodes (83.7\%, Wilson 95\% CI
$[79.1,87.4]$), compared with 10/300 for the equal-parameter Geometry model
(3.3\%, $[1.8,6.0]$).  The paired success differences at difficulties 0.2,
0.5, and 0.8 are 0.90, 0.83, and 0.68; the corresponding Holm-adjusted exact
McNemar values are $5.33\times10^{-26}$, $5.79\times10^{-24}$, and
$2.29\times10^{-18}$.  Across all difficulties, 242 physical cells succeed only
with TopoHarness, one succeeds only with Geometry, nine succeed with both, and
48 fail with both.  The scripted teacher succeeds in 300/300 episodes, showing
that the residual 49 TopoHarness failures are policy errors rather than an
infeasible test set.

The architecture ablation changes the interpretation.  ACT-RelPool,
Semantic-only, the no-unary-ID/order-feature model, and the 20\% and 50\% data
models all produce exactly the same 251/300 success count and failure partition
as TopoHarness.  Physical-only reaches 1/300, direct ACT and ACT-Pointer reach
0/300, and random control reaches 0/300.  Thus this environment identifies the
necessity of an observable process relation and a relation-grounded action
interface, but it does not identify learned message passing, unary roles, or the
full training-set size as separately necessary.  We treat this saturation as a
benchmark limitation rather than an architectural win.

\begin{table}[t]
  \centering
  \caption{Direct MuJoCo pose-observation/task-space-control audit.  The
  detailed robot is a visual IK twin, not a torque-controlled arm.}
  \label{tab:mujoco-direct}
  \IfFileExists{../../artifacts/papers/routebot/final/mujoco_direct/route_mujoco_direct_table.tex}
    {\input{../../artifacts/papers/routebot/final/mujoco_direct/route_mujoco_direct_table.tex}}
    {\fbox{\parbox{\dimexpr\columnwidth-2\fboxsep-2\fboxrule\relax}{%
      \centering Frozen table generated by
      \texttt{sim-eval-route-mujoco-direct}}}}
\end{table}

\paragraph{Frozen direct MuJoCo audit.}
Across 20 paired fixture perturbations and all nine segment assignments,
TopoHarness and the scripted teacher each complete 180/180 episodes; the
equal-parameter geometry policy completes 0/180.  The Wilson interval is
$[0.979,1.000]$ for TopoHarness and $[0.000,0.021]$ for Geometry.  All 180
discordant pairs favor TopoHarness (two-sided exact McNemar test with
Holm adjustment, $p=2.61\times10^{-54}$).  This isolates relation-sensitive
closed-loop decisions under a second cable dynamics backend, but the mocap
control abstraction prevents a joint-level or force-control claim.

\begin{figure}[t]
  \centering
  \includegraphics[width=\columnwidth]{%
    ../../artifacts/papers/routebot/final/mujoco_direct_visual/route_mujoco_direct_final.png}
  \caption{Frozen TopoHarness final state under direct MuJoCo observation and
  task-space action control (assignment 7--15--23).  The three semantic
  latches and endpoint succeed; the rendered UR10e is a visual IK twin.}
  \label{fig:mujoco-direct}
\end{figure}

\begin{table}[t]
  \centering
  \caption{Episode-disjoint offline action prediction on the public real-robot
  Berkeley Cable Routing data.  These are not closed-loop success rates.}
  \label{tab:public-real}
  \IfFileExists{../../artifacts/papers/routebot/public_real/route_public_table.tex}
    {{\small\setlength{\tabcolsep}{1.5pt}%
      \input{../../artifacts/papers/routebot/public_real/route_public_table.tex}}}
    {\fbox{\parbox{\dimexpr\columnwidth-2\fboxsep-2\fboxrule\relax}{%
      \centering Frozen public real-data audit table}}}
\end{table}

\paragraph{Frozen public real-data audit.}
On 247 held-out real-robot episodes, adding the previously executed action to
the state ridge reduces mean per-episode active-action RMSE by 11.6\% relative
to action persistence (paired difference $-0.0191$, 95\% episode-bootstrap CI
$[-0.0206,-0.0176]$).  It improves 233 of 247 episodes (two-sided exact sign
test $p=2.33\times10^{-52}$).  This confirms that the data pipeline detects a
real temporal signal; it cannot test semantic binding or closed-loop routing.

\begin{table}[!b]
  \centering
  \caption{Exact semantic counterfactual over nine assignments. Each policy
  observes identical physical states within a seed; only relation edges and the
  corresponding success predicate change.}
  \label{tab:counterfactual}
  \IfFileExists{../../artifacts/papers/routebot/final/counterfactual/route_counterfactual_table.tex}
    {{\small\setlength{\tabcolsep}{2.5pt}%
      \input{../../artifacts/papers/routebot/final/counterfactual/route_counterfactual_table.tex}}}
    {\fbox{\parbox{\dimexpr\columnwidth-2\fboxsep-2\fboxrule\relax}{%
      \centering Frozen table generated by
      \texttt{sim-eval-route-counterfactual}}}}
\end{table}
\FloatBarrier

\paragraph{Frozen nine-assignment counterfactual.}
Across 2,700 physical-seed--assignment cells per policy, TopoHarness succeeds
in 2,340 (86.7\%, Wilson 95\% CI $[85.3,87.9]$), versus 57 for Geometry
(2.1\%, $[1.6,2.7]$), 268 for ACT-Pointer (9.9\%, $[8.9,11.1]$), and 2,700
for the teacher.  TopoHarness alone succeeds in 2,284 Geometry-paired cells;
Geometry alone succeeds in one, while 56 succeed and 359 fail with both.  The
aggregate exact McNemar value is $6.41\times10^{-685}$.  More conservatively,
each of the 27 assignment--difficulty comparisons has a 60--100 percentage
point advantage over Geometry, and the largest Holm-adjusted value is
$1.06\times10^{-16}$.  ACT-RelPool again has exactly the same 2,340 successes
and failure partition as TopoHarness.  Its per-assignment success spans
83--100\%, 76--100\%, and 60--94\% at increasing difficulty; the wider range
at high difficulty bounds any claim of perfect assignment invariance.

\begin{figure}[!b]
  \centering
  \includegraphics[width=\columnwidth]{%
    ../../artifacts/papers/routebot/final/counterfactual/route_counterfactual_heatmap.pdf}
  \caption{Cellwise semantic success for the exact relation intervention
  (100 paired physical seeds per cell; shared 0--100\% scale). TopoHarness and
  ACT-RelPool are cellwise identical. ACT-Pointer succeeds only on isolated
  assignments, revealing assignment-specific shortcuts.}
  \label{fig:counterfactual-heatmap}
\end{figure}

\begin{table}[t]
  \centering
  \caption{Matched open-loop sanity check in official DLO-Lab
  \texttt{wiring\_post}. Native $R$ is unmodified; relation and winding
  diagnostics are RouteBot extensions.}
  \label{tab:dlolab}
  \IfFileExists{../../artifacts/papers/routebot/final/dlolab_external/matched_paths/matched_path_table.tex}
    {\resizebox{\columnwidth}{!}{%
      \input{../../artifacts/papers/routebot/final/dlolab_external/matched_paths/matched_path_table.tex}}}
    {\fbox{\parbox{\dimexpr\columnwidth-2\fboxsep-2\fboxrule\relax}{%
      \centering Frozen DLO-Lab matched-path table}}}
\end{table}
\FloatBarrier

\paragraph{Official DLO-Lab execution sanity check.}
All four 30-vertex rollouts remain finite. The ordered fixture route reaches
native reward 0.975 and ordered-point RMSE 1.46 cm, versus 0.571 and 12.30 cm
for the wrong-order path. Its fixture-relation angular error is 0.073 rad,
compared with 0.785 rad for wrong order. However, the endpoint-straight path
also reaches reward 0.970, 2.73 cm RMSE, and a 0.069-rad relation error.
Therefore this result validates the pinned external environment and detects a
gross wrong-order path, but the official geometric target does not isolate
RouteBot's semantic intervention. We do not count it as trained external
semantic replication.

\section{Limitations and Claim Boundary}

The current environments are simulated, the PBD backend is two-dimensional and
uncalibrated, and the detailed robot arms do not participate in cable contact.
No result in this paper is a production-line success rate.  The counterfactual
isolates a process relation within a simulator; it does not establish causal
effects on real hardware. Force-controlled clip insertion, perception from
raw images, cable-parameter identification, and trained external closed-loop
semantic replication remain mandatory before a hardware-level claim. The
single-seed DLO-Lab matched-path check is open loop and its endpoint-straight
control nearly matches the ordered route, so it cannot satisfy that semantic
replication requirement. The direct MuJoCo audit uses
current three-dimensional object poses but a mocap task-space gripper, so it
does not validate contact-force regulation or joint-level control.  The
public-data experiment is an offline action audit and does not satisfy the
external replication gate.

\section{Ethical and Societal Considerations}

The study uses synthetic simulator trajectories and attributed public robot
teleoperation traces; we conduct no new human-subject data collection and use
no personal data.  Wire-harness
automation could reduce repetitive ergonomic strain, but it
could also displace manual assembly work or intensify monitoring if deployed
without worker participation.  The present controller is not safety certified:
its outputs must not command production robots without guarded workcells, risk
assessment, force limits, and qualified human supervision.  We therefore report
only simulator evidence and publish failure cases as well as successes.  The
models are small CPU-trainable baselines; training time and hardware are included
in the reproducibility record so computational cost is auditable.

\section{Reproducibility}

The anonymous package will include episode-disjoint demonstrations, hashes,
frozen checkpoints, all seeds and configurations, per-episode CSV files,
statistical scripts, MJCF sources, retained third-party licenses, and commands
that regenerate every table and figure.  The independent ACT adaptation states
its differences from the official implementation and does not copy upstream
source code. The external DLO-Lab asset archive is not redistributed because
it contains no asset license; the package records its official URL, exact hash,
required-file hash, and a deterministic local verification command.

\balance
\bibliographystyle{plain}
\bibliography{references}

\end{document}
